\begin{abcliste}
\abc With $\dot I = \ddot Q$ the equation of the circuit reads
\begin{align}
 L\,\ddot Q &= -\frac{1}{C}\,Q
\end{align}
Term by term with $m\,\ddot y = -D\,y$: the charge $Q$ plays the role of $\result{y}$, the current $I = \dot Q$ that of $\result{v = \dot y}$, the inductance $L$ that of the $\result{\text{mass } m}$, and $1/C$ that of the $\result{\text{spring constant } D}$ (a large capacitance is like a soft spring).

\abc As $\omega = \sqrt{D/m}$ for the spring:
\begin{align}
 \omega &= \omF = \frac{1}{\sqrt{\Li\times\Cp}} = \om \approx \result{\omP{0}{3}} \\
 T &= \TF = \T \approx \result{\TP{-3}{3}}
\end{align}

\abc At $t = 0$ the capacitor is fully charged and no current flows yet, like a body let go at its turning point. The current is largest when the capacitor is empty for the first time, like the body passing the equilibrium, a quarter of a period later:
\begin{align}
 t &= \tqF = \tq \approx \result{\tqP{-3}{3}}
\end{align}
The charge oscillates as $Q = Q_0\cos(\omega t)$ with $Q_0 = C\,U_0 = \QzP{-6}{3}$; like $\hat v = \omega A$, the largest current is
\begin{align}
 \hat I &= \IxF = \om\times\Qz = \Ix \approx \result{\IxP{-3}{3}}
\end{align}
\end{abcliste}
